% =====================================================================
%  config/comandos.tex — comandos auxiliares
%
%  Citações no estilo da SBC:
%    \cite{alongi2022}                       -> [Alongi et al. 2022]
%    \cite{alongi2022,borges2024}            -> [Alongi et al. 2022, Borges et al. 2024]
%    Alongi et al. \citeano{alongi2022}      -> Alongi et al. [2022]  (citação narrativa)
% =====================================================================
\makeatletter
\newcommand{\citeano}[1]{%
  \nocite{#1}%
  \@ifundefined{b@#1}%
    {[\textbf{?}]\@latex@warning{Citation `#1' undefined}}%
    {\begingroup
       % com hyperref, \b@chave guarda \hyper@@link[cite]{}{cite.chave}{rótulo}
       \def\hyper@@link[##1]##2##3##4{##4}%
       \protected@xdef\citeano@rotulo{\csname b@#1\endcsname}%
     \endgroup
     [\hyperlink{cite.#1}{\expandafter\citeano@ano\citeano@rotulo\@nil}]}}
% o ano é a última palavra do rótulo ("Amaro et~al. 2026" -> 2026)
\def\citeano@ano#1\@nil{\citeano@ultima#1 \@nil}
\def\citeano@ultima#1 #2\@nil{%
  \if\relax\detokenize{#2}\relax
    \expandafter\@firstoftwo
  \else
    \expandafter\@secondoftwo
  \fi
  {#1}{\citeano@ultima#2\@nil}}
\makeatother

% O sbc-template.sty (v2017) escreve "Plavras-Chave"; corrigido aqui.
\renewcommand{\palavraschave}[1]{%
  \vspace{6pt}\small\linebreak\textbf{\textit{Palavras-chave---}} #1}
